
\section{Conclusions}
\label{sec:conclusions}

This study developed a hierarchical aerodynamic knowledge-transfer framework for data-efficient prediction of three-dimensional surface-pressure distributions on transonic swept wings. Reusable sectional pressure knowledge learned from infinite-span swept-wing data and localized wing-root and wing-tip corrections derived from paired finite- and infinite-span CFD data were mapped to the target configuration to construct a hierarchical pressure prior. A target-domain residual model was then trained using limited high-fidelity CFD samples to learn the remaining pressure discrepancy.

The results demonstrate the effectiveness of hierarchical knowledge transfer for data-efficient three-dimensional pressure prediction. With only 10 target-domain training wings, M2 reduced the surface-pressure RMSE by 65.74\% relative to the direct-learning baseline M0 and maintained the lowest mean prediction error across all investigated training budgets. Regional ablation results further demonstrate the complementary contributions of the root and tip knowledge models. At $N=20$, M2 reduced the mean RMSE relative to M1 by 39.81\%, 24.69\%, and 20.66\% in the root, midspan, and tip regions, respectively. These improvements indicate that regional knowledge provides benefits beyond the immediate vicinity of the corresponding spanwise boundaries, although the magnitude of improvement varies with spanwise position and surface.

The predicted pressure distributions also provide useful information for downstream aerodynamic-coefficient regression under limited force-label supervision. With 20 force-labeled target wings, D1, which uses pressure predictions from the fixed M2 model, reduced the mean absolute errors in lift, drag, and pitching-moment coefficients by 70.1\%, 54.8\%, and 23.4\%, respectively, relative to direct regression D0. The improvements were more pronounced for lift and drag than for pitching moment, indicating that the predictive benefit of pressure information depends on the aerodynamic quantity being modeled.

The principal contribution of this work lies in developing a hierarchical representation and transfer strategy that separates the reusable sectional aerodynamic response from localized three-dimensional pressure corrections. These complementary knowledge components are mapped and combined to construct a pressure prior for the target wing. This formulation shifts target-domain learning from directly fitting the entire three-dimensional pressure field to correcting the pressure variations not adequately captured by the transferred knowledge. Consequently, limited target-domain CFD samples can be more effectively utilized to learn configuration-dependent three-dimensional pressure characteristics rather than repeatedly learning the sectional responses already represented by the prior. This approach provides a basis for exploiting hierarchical aerodynamic knowledge in small-sample wing pressure prediction and for further integrating distributed pressure information into aerodynamic performance modeling.
